
\subsection{Computing the extensions}\label{sec:extensions}

Between the $\mathrm{E}^2= \mathrm{E}^\infty$-page of the one step spectral sequence
\[
  \THH_*(\BP \langle 1 \rangle)\langle \sigma v_2 \rangle \Rightarrow \THH_*(\BP \langle 2 \rangle ; \BP \langle 1 \rangle)
\]
and the target module $\THH_*(\BP \langle 2 \rangle ; \BP \langle 1 \rangle)$, we have to solve extension problems of the following form: an element $\alpha$ of the $\mathrm{E}^\infty$-page, non divisible by $\sigma v_2$ and such that $p \cdot \alpha = 0$, can lift to an element $\bar{\alpha} \in \THH_*(\BP \langle 2 \rangle ; \BP \langle 1 \rangle)$ such that $p \cdot \bar{\alpha} = \beta$ where $\beta$ is a lift of a multiple of $\sigma v_2$. We will solve these extension problems using what we know of $\THH_*(\BP \langle 2 \rangle ; \Z_{(p)})$. 

\begin{lemma}\label{lem432}
  There can be no $\mathrm{BP}\langle 1\rangle_*$-module extensions between torsion elements.
\end{lemma}

\begin{proof}
  The torsion elements are in degrees
  \[
    \begin{gathered}
      |v_1^k b_i | = 2k(p-1) + 2i p^2 + 2p - 2 \\
      |v_1^h b_j \sigma v_2| = 2h(p-1) + 2(j+1)p^2 + 2p - 3
    \end{gathered}
  \]
  respectively for the non mutliple and multiple of $\sigma v_2$. Since the former is even and the latter odd, there can be no $\mathrm{BP}\langle 1\rangle_*$-module extensions.
\end{proof}


\begin{lemma} \label{lem:onlyextensions}
  The only possible $\mathrm{BP}\langle 1\rangle_*$-module extensions %
  are of the form
  \[
    p \cdot b_1 = \lambda_1 \sigma v_2
  \]
  and
  \[
    p \cdot v_0^k b_{p^k} = v_1^{p^{k+1} - p}v_0^{k-1} a_{p^{k-1}} \sigma v_2
  \]
  for $k \geq 1$.
\end{lemma}

\begin{proof}
  We have to investigate the extensions between a non $\sigma v_2$-multiple of  a torsion element and a non-torsion $\sigma v_2$-multiple element. As in Lemma~\ref{lem432}, the non-torsion element needs to be in even degree, that is of the form
  \begin{equation} \label{eq:torsionliftform}
    v_1^h p^k \lambda_1 \sigma v_2 \text{ or } v_1^h p^k v_0^n a_{p^n} \sigma v_2.
  \end{equation}

  Remember the spectral sequence \eqref{BrunSS} is a long exact sequence of the form:
  \[
    \dots \rightarrow \THH_*(\BP \langle 1 \rangle)\{ \sigma v_2\} \rightarrow \THH_*(\BP \langle 2 \rangle ; \BP \langle 1 \rangle) \rightarrow \THH_*(\BP \langle 1 \rangle)\{1\} \rightarrow \dots
  \]
  Let $\alpha \in \THH_*(\BP \langle 1 \rangle)\{1\}$ be a torsion element in the $E^\infty$ page, such that $p \cdot \alpha = 0$. Write $m$ for the smallest integer such that $v_1^m \alpha = 0$. Let $\beta \sigma v_2 \in \THH_*(\BP \langle 1 \rangle)\{ \sigma v_2\}$ be of the form (\ref{eq:torsionliftform}), that is of a non torsion element of even degree. Write $\bar{\beta} \sigma v_2$ for its image in $\THH_*(\BP \langle 2 \rangle ; \BP \langle 1 \rangle)$. Suppose that $\beta \sigma v_2$ is divisible by $p$, with $p \cdot \beta' \sigma v_2 = \beta \sigma v_2$, and write $\bar{\beta'} \sigma v_2$ for its images, so that $p \cdot \bar{\beta'} \sigma v_2 = \bar{\beta} \sigma v_2$.

  If the extension $p \cdot \bar{\alpha} = \bar{\beta} \sigma v_2$ occurs, in $\THH_*(\BP \langle 2 \rangle ; \BP \langle 1 \rangle)$ we must have
  \[
    p v_1^m \bar{\alpha} = v_1^m p \bar{\alpha} = v_1^m \beta \sigma v_2 \neq 0
  \]
  so that $v_1^m \bar{\alpha} \neq 0$. However, the image of $v_1^m \bar{\alpha}$ in $\THH_*(\BP \langle 1 \rangle)\{1\}$ is $v_1^m \alpha = 0$, so that $v_1^m \bar{\alpha}$ must lift to $\THH_*(\BP \langle 1 \rangle)\{ \sigma v_2\}$. Since the non-torsion generators in even degree in $\THH_*(\BP \langle 1 \rangle)\{ \sigma v_2\}$ are unique up to a unit, the only possible lift is $v_1^m \beta' \sigma v_2$, which implies that in $\THH_*(\BP \langle 2 \rangle ; \BP \langle 1 \rangle)$ we have an equality
  \[
    v_1^m \bar{\alpha} = v_1^m \bar{\beta'} \sigma v_2.
  \]
  Let
  \[
    \hat{\alpha} = \bar{\alpha} - \bar{\beta'} \sigma v_2
  \]
  be a new lift for $\alpha$. We have $p \cdot \hat{\alpha} = 0$ and $v_1^m \hat{\alpha} = 0$, so that we have found a lift that has the same properties as $\alpha \in \mathrm{E}^\infty$.

  Thus, the only extensions we need to care about are the ones where $\beta \sigma v_2$ is not already divisible by $p$. However, it must still be the case that $v_1^m \bar{\alpha} \neq 0$, so that $v_1^m \beta \sigma v_2$ must be divisible by $p$ in $\THH_*(\BP \langle 1 \rangle)\{\sigma v_2\}$.

  Recall the relations in the non-torsion of $\THH_*(\BP \langle 1 \rangle)$
  \[
    \begin{gathered}
      p\cdot a_1 = v_1^p \lambda_1 \\
      p\cdot v_0^n a_{p^n} = v_1^{p^{n+1}}v_0^{n-1} a_{p^{n-1}}
    \end{gathered}
  \]
  whose degrees in $\THH_*(\BP \langle 1 \rangle)\{ \sigma v_2\}$ are $|v_0^n a_{p^n}\sigma v_2| = 2(p^n + 1)p^2 - 2$. The only suitable $\alpha$ in the $\mathrm{E}^\infty$-page with $v_1^m \alpha$ of the right degree are the element of $M$, that is the $v_0^k b_{p^k}$, since in general we have
  \[
    |v_1^{p^{\nu_p(i) - h + 1 } + \dots + p} v_0^h b_i| = 2(i + p^{\nu_p(i) - h})p^2 - 2 \,.
  \]
\end{proof}

\begin{proposition}
  The $\BP\langle 1\rangle_*$-module extensions of Lemma \ref{lem:onlyextensions} all occur. 
\end{proposition}

\begin{proof}
  If these extensions do not occur, we have elements $v_0^k b_{p^k}$ in $\THH_*(\BP \langle 2 \rangle ; \BP \langle 1 \rangle)$ such that
  \[
    v_1^p v_0^k b_{p^k} = 0
  \]
  so that there must be a $v_1$-Bockstein in the spectral sequence
  \[
    \THH_*(\BP \langle 2 \rangle ; \Z_{(p)})[v_1] \Rightarrow \THH_*(\BP \langle 2 \rangle ; \BP \langle 1 \rangle)
  \]
  that is a differential whose source must be in degree
  \[
    |v_1^p v_0^k b_{p^k}| + 1 = 2(p^k + 1)p^2 - 1.
  \]
  However, by Theorem~\ref{thm:BP2-coeff} there is no torsion element in $\THH_*(\BP \langle 2 \rangle; \Z_{(p)})$ in that degree, since they are in degrees
  \[
    2ip^3 - 1 ,\, 2ip^3 - 1 + 2p - 1,\, 2ip^3 - 1 + 2p^2 - 1 \text{ and } 2ip^3 - 1 + 2p^2 - 1 + 2p - 1 \,.
  \]
\end{proof}

We are now prepared to prove our main theorem.


\begin{theorem}\label{thm:main-v2}
Suppose $\BP\langle 2\rangle$ satisfies Running Assumption~\ref{runningassumption}. 
The topological Hochschild homology of $\BP\langle 2\rangle$ is a direct sum
  \[
    \THH_*(\BP \langle 2 \rangle ; \BP \langle 1 \rangle) \cong \Z_{(p)}[v_1] \langle \sigma v_2 \rangle \oplus \mathcal{F} \oplus \Sigma^{2p - 1} (\mathcal{F}_{\geq 2p^2 - 1}) \oplus \mathcal{T} \oplus \Sigma^{2p - 1} \mathcal{T}
  \]
  where $\mathcal{T}$ is the $\Z_{(p)}[v_1]$-submodule of $\THH_*(\BP \langle 1 \rangle)$ containing the torsion elements that are in the image of the multiplication by $v_1^p$ map,  $\mathcal{F}$ denotes the submodule of the torsion free part of $\THH_*(\BP \langle 1 \rangle)$ of elements not of the form $v_1^k \cdot 1$ for some $k \geq 0$, and $\mathcal{F}_{\geq 2p^2 -1}$ is the submodule of $\mathcal{F}$ whose elements are in degree $\geq 2p^2 -1$.
\end{theorem}

\begin{proof}[Proof of Theorem~\ref{thm:main-v2}]
Let $\mathcal{T}$ denote the $\Z_{(p)}[v_1]$-submodule of $\THH_*(\BP \langle 1 \rangle)$ containing the torsion elements that are in the image of the multiplication by $v_1^p$ map. Then $\mathcal{T}$ corresponds to the kernel of the differential on the torsion as described in Proposition \ref{prop:specseqtorsion}. In the torsion of $\THH_*(\BP \langle 1 \rangle)$, there is an isomorphism of $\Z_{(p)}[v_1]$-modules
  \begin{center}
    \begin{tikzcd}[row sep=tiny]
      \im (\times v_1^p) \rar & \operatorname{coim}(\times v_1^p) \\
      v_1^p b_i \rar[mapsto] & b_i
    \end{tikzcd}
  \end{center}
  that lowers the degree by $|v_1^p| = 2p^2 - 2p$. But in the spectral sequence, the second copy $\THH_*(\BP \langle 1 \rangle)\{ \sigma v_2\}$ is in degree $|\sigma v_2| = 2p^2 - 1$ higher, so in the end the shift is by $2p -1$. The torsion from the $\sigma v_2$-copy therefore agrees with $\Sigma^{2p-1} \mathcal{T}$.

  The $v_1$-towers above $1$ and $\sigma v_2$ are not touched by extensions, so that we get the $\Z_{(p)}[v_1]\langle \sigma v_2 \rangle$ summand. The description of $\mathcal{F}$ in the statement is equivalent to have
  \[
    \text{torsion free part of } \THH_*(\BP \langle 1 \rangle) \cong \Z_{(p)}[v_1] \{1\} \oplus \mathcal{F}
  \]
  so the $\mathcal{F}$ in the result comes directly from the torsion free part of the corresponding $\THH_*(\BP \langle 1 \rangle)\{1\}$ copy in the spectral sequence.

  We still have to account for the elements of the submodule $M$ of Proposition \ref{prop:specseqtorsion}. They form extensions with non torsion elements of $\THH_*(\BP \langle 1 \rangle)\{ \sigma v_2\}$, which are all permanent cycles and not boundaries. The remaining torsion free classes coming from $\THH_*(\BP \langle 1 \rangle)\{ \sigma v_2\}$ and the classes in $M$ form a $\Z_{(p)}[v_1]$-module isomorphic to $\Sigma^{2p-1}(\mathcal{F}_{\geq 2p^2 - 1})$ with the isomorphism given by the map
  \[
    v_0^n a_{p^n} \mapsto v_0^n b_{p^n},\, n \geq 0
  \]
  where the $v_0^n a_{p^n}$ are the $\Z_{(p)}[v_1]$ generators of $\mathcal{F}_{\geq 2p^2 - 1}$.
\end{proof}


\begin{remark}
  We can derive an explicit description of $\mathcal{T}$  and $\mathcal{F}$ from $\THH_*(\BP \langle 1 \rangle)$. We write $c_i$ for the lift of the class $v_1^p b_i$ of the spectral sequence.

  $\mathcal{T}$ is a quotient of the $\Z_{(p)}[v_1]$-module
  \[
    \Z_{(p)}[v_1]\{v_0^h c_{\alpha p^n
    },\, n\geq 1,\, \alpha\geq 1,\, p\nmid \alpha,\, h \geq 0\}
  \]
  by the relations:
  \begin{itemize}
  \item $v_0^{h} c_{\alpha p^n} = 0$ for any $\alpha \geq 1$ and $n \geq 1$, $\alpha$ not divisible by $p$, and $h \geq n$,
  \item $v_1^{p^{n-h+1}+p^{n-h}+\dots+p^2 }\cdot v_0^h c_{\alpha p^n} = 0$ for any $\alpha \geq 1$ and $n \geq 1$, $\alpha$ not divisible by $p$ and  $0 \leq h \leq n - 1$,
  \item $p\cdot c_{(\beta p+p-1)p^n} = v_0 c_{(\beta p+p-1)p^n}+v_1^{p^{n+2}}v_0^{\nu_p(\beta)} c_{\beta p^{n+1}}$ for any $\beta \geq 1$ and $n \geq 1$.
  \item $p\cdot v_0^h c_{\alpha p^n} = v_0^{h+1} c_{\alpha p^n}$ for any $\alpha \geq 1$, $n \geq 1$, $\alpha$ not divisible by $p$, and any $1 \leq h \leq n$, or $h=0$ not in the previous case.
  \end{itemize}

  $\mathcal{F}$ is a quotient of the $\Z_{(p)}[v_1]$-module
  \[
    \Z_{(p)}[v_1]\{\lambda_1,\, v_0^n a_{p^n},\,n\geq 0\} 
  \]
    by the relations:
  \begin{itemize}
  \item $p\cdot a_1 = v_1^p \lambda_1$,
  \item $p\cdot v_0^n a_{p^n} = v_1^{p^{n+1}}v_0^{n-1} a_{p^{n-1}}$ for any $n\geq 1$,
  \end{itemize}
  $\mathcal{F}_{\geq 2p^2 -1}$ is obtained from $\mathcal{F}$ by removing the generator $\lambda_1$.

  $\nu_p$ is the $p$-adic valuation, and the degrees are
  \[
    \begin{gathered}
      |a_i| = 2ip^2 - 1 \,, \\
      |c_i| = 2ip^2 - 1 + 2p^2 - 1 \,.
    \end{gathered}
  \]
\end{remark}


\begin{remark}
  We can describe
  \[
    \THH_*(\BP \langle 1 \rangle; \BP \langle 0 \rangle) \cong \THH_*(\ell; \Z_{(p)}) \cong \Z_{(p)}\langle \sigma v_1 \rangle \oplus \mathcal{T}_1 \oplus \Sigma^{2p - 1} \mathcal{T}_1
  \]
  where $\mathcal{T}_1$ consists of the torsion elements of $\THH_*(\BP \langle 0 \rangle)$ divisible by $p = v_0^1$.

  We have just described (we add an index here to disambiguate)
  \[
    \text{torsion of } \THH_*(\BP \langle 2 \rangle; \BP \langle 1 \rangle)\cong  \mathcal{T}_2 \oplus \Sigma^{2p - 1} \mathcal{T}_2
  \]
  where $\mathcal{T}_2$ consists the torsion elements of $\THH_*(\BP \langle 1 \rangle)$ divisible by $v_1^p$.

  Going down on more height, we can describe
  \[
    \THH_*(\BP \langle 0 \rangle; \BP \langle -1 \rangle) \cong \THH_*(H \Z_{(p)}; H\F_p) \cong \mathcal{T}_0 \oplus \Sigma^{2p - 1} \mathcal{T}_0
  \]
  where $\mathcal{T}_0$ is the image of the $p$-th power map. Here there is no torsion free part. 
  
  We speculate that for all $n \geq 1$, we have
  \[
    \text{torsion of } \THH_*(\BP \langle n \rangle; \BP \langle n-1 \rangle) \cong  \mathcal{T}_n \oplus \Sigma^{2p - 1} \mathcal{T}_n
  \]
  where $\mathcal{T}_n$ is the torsion elements of $\THH_*(\BP \langle n-1 \rangle)$ divisible by $v_{n-1}^{p^{n-1}}$.

  %

  %
  %
  
  Moreover, the isomorphism between $\mathcal{T}_1$ and $\Sigma^{2p-1} \mathcal{T}_1$ is obtained by some operation  denoted $v_0^{-1} \sigma v_1$, the isomorphism between $\mathcal{T}_2$ and $\Sigma^{2p-1} \mathcal{T}_2$ is obtained by some operation denoted $v_1^{-p} \sigma v_2$, so that the isomorphism between $\mathcal{T}_n$ and $\Sigma^{2p-1} \mathcal{T}_n$ could be obtained by the operation $v_{n-1}^{-p^{n-1}} \sigma v_n$ in the Brun spectral sequence of Proposition~\ref{prop:brunss}.
\end{remark}


\section{Truncated Brown--Peterson spectra are not Thom spectra}\label{notThom}

In this section, we implicitly $2$-localize and write $\ku=\ell=\BP\langle 1\rangle$ and $\BP\langle n\rangle$ denotes an $\mathbb{E}_3$-$\MU$-algebra form of the $n$-th truncated Brown--Peterson spectrum at the prime $p=2$. 
The goal of this section is to prove Theorem~\ref{thm:main-thom}. 

First we need some notation. Given an $\mathbb{E}_2$-ring $A$ and an $\mathbb{E}_1$-$R$-algebra $B$ , we write 
 $\overline{\mathrm{THH}}(A;B)$ for the cofiber of the unit map 
$B\to \mathrm{THH}(A;B)$ in $B$-modules. The only input for our proof is the following computational result. 

\begin{proposition}\label{prop-computational-input}
We compute that $\pi_*\overline{\mathrm{THH}}(\BP\langle n\rangle;\ku)$  is generated as a $ku_*$-module in degrees $\le 8$ by $\lambda_1$, $a_1$ and $\sigma v_2$ in degrees $3$, $7$, $7$ respectively and there is a relation $pa_1= v_1^2 \lambda_1$
\end{proposition}

\begin{proof}
We consider the fiber sequence  
\[ \THH(\ku;F) \to  \THH(\BP\langle n\rangle ; \ku )\to 
\THH(\ku) 
\]
where $F$ is the fiber of the map 
\[ \ku \otimes_{\BP\langle n \rangle }\ku \to \ku \]
of $\ku\otimes \ku^{\textup{op}}$-modules. 
The map $\pi_*\ku\to \pi_*\THH(\ku)$ is an isomorphism in degrees $*\le 2$ by Theorem~\ref{thm:AHL}. Since we also know that  
\[ 
\pi_k F = \begin{cases} 0 & \text{ if } k<7\text{ or } k = 8,\\
           \mathbb{Z}_{(p)} &\text{ if } k=7  
\end{cases}
\]
we can conclude that the inclusion of $0$-simplices $F\to \THH(\ku;F)$ is an isomorphism in degrees $\le 8$. We write $\sigma v_2$ for the image of the generator of $\THH_{7}(\ku;F)=\mathbb{Z}_{(p)}$ in $\THH_7(\BP\langle n\rangle ;\ku)$. Note that $\THH_7(\ku)=\mathbb{Z}_{(p)}\{a_1\}$ so the extension is trivial and we have proven the claim. 
\end{proof}

\begin{proof}[Proof of Theorem~\ref{thm:main-thom}] 
Suppose $\BP\langle n\rangle$ is the Thom spectrum of a $2$-fold loop map at the prime $2$. Then 
\[ \THH(\BP\langle n\rangle;\mathrm{ku})\simeq \mathrm{ku} \otimes X_{+}\simeq \mathrm{ku} \otimes_{\mathrm{ko}}(\mathrm{ko} \otimes X_+)
\]
by~\cite[Theorem~2]{BCS10} together with the fact that $\eta$ has trivial image under the Hurewicz map $\mathbb{S}\to \BP\langle n\rangle$, see the proof of~\cite[Theorem~2]{Bea17}. We argue that no such $Y=\mathrm{ko}\otimes X_{+}$ exists by the same argument as~\cite{AHL09}. This relies on Proposition~\ref{prop-computational-input}. We include the details for completeness since the argument from~\cite{AHL09} needs to be adapted slightly. 

Suppose a cellular $ko$-module $Y$ existed. If it did, it would have cells beginning in degrees $3$, $7$, $7$ and $8$ by Proposition~\ref{prop-computational-input}. Since $\pi_3\mathrm{ku}=0$, we know the first attaching map of $X$ is zero and therefore the first attaching maps of $Y$ is also trivial. We therefore know that $Y$ begins with 
\[ \Sigma^{3}\mathrm{ko}\vee \Sigma^{7}\mathrm{ko}\vee \Sigma^{7}\mathrm{ko}\to (\Sigma^{3}\mathrm{ko}\vee \Sigma^{7}\mathrm{ko}\vee \Sigma^{7}\mathrm{ko})\cup_{\phi} C\Sigma^{7}\mathrm{ko})\to Y
\]
The attaching map $\phi$ is an element in 
\begin{equation}\label{pi7ko}\pi_{7} (\Sigma^{3}\mathrm{ko}\vee \Sigma^{7}\mathrm{ko}\vee \Sigma^{7}\mathrm{ko})=\mathbb{Z}\{v_1^2\}\oplus \mathbb{Z}\{1\}\oplus \mathbb{Z}\{1\} \,.
\end{equation}
By the relation $v_1^2\lambda_1=a_1$, the goal would be to consider the class 
\[
(v_1^2,1,0) \in \pi_{7} (\Sigma^{3}\mathrm{ko}\vee \Sigma^{7}\mathrm{ko}\vee \Sigma^{7}\mathrm{ko})=\mathbb{Z}\{v_1^2\}\oplus \mathbb{Z}\{1\}\oplus \mathbb{Z}\{1\} 
\]
and lift it to \eqref{pi7ko}. However the map 
\[
\pi_7(\Sigma^{3}\mathrm{ko}\vee \Sigma^{7}\mathrm{ko}\vee \Sigma^{7}\mathrm{ko})  \to 
\pi_7(\Sigma^{3}\mathrm{ku}\vee \Sigma^{7}\mathrm{ku}\vee \Sigma^{7}\mathrm{ku}) 
\]
is induced by the complexification map $\mathrm{ko}\to \mathrm{ku}$, so it is multiplication by 2 on the $\Sigma^3$ summand and identity on the $\Sigma^7$ summands, and therefore $(v_1^2,1,0)$ is not in the image, leading to a contradiction. 
\end{proof}

\begin{remark}
As in the case of~\cite{AHL09}, our argument for proving Theorem~\ref{thm:main-thom} requires that we work at the prime $p=2$. 
\end{remark}

\begin{remark}
We have proven that $\BP\langle n\rangle$ at $p=2$ is not a Thom spectrum of a $2$-fold loop map over the sphere spectrum, however it may still be a Thom spectrum in the category of modules over a different spectrum in the sense of~\cite{ABGHR14}. For example, it has been conjectured that it is the Thom spectrum over certain Ravenel spectra~\cite[Conjecture~1.1.3]{Dev24} and it has been proven that it is the Thom spectrum in the category of modules over a polynomial algebra over complex cobordism in~\cite{Qui25}. 
\end{remark}
